% license-footer.tex -- shared LaTeX preamble that stamps the open-educational
% material branding ("CC BY 4.0 · © 2026 Robert Ioffe") as a page footer.  Included
% via pandoc's --include-in-header from tools/build-book-localtoc.sh.
%
% IMPLEMENTATION NOTE (why centered, not per-side):
% This PDF is a two-sided `book` class rendered through
% pandoc + xelatex + hyperref (hyperref is loaded by the two-level-TOC preamble).
% In that combination fancyhdr's per-side position codes (LO/RO/LE/RE) are
% unreliable -- they can drop the footer on EVEN pages.  A single centered footer
% with an \hfill ("license <tab> pagenum") is robust: it lands on every page, odd
% and even alike, and \hfill still flushes the license to the left and the page
% number to the right.  We reapply it to every page style a book uses -- "fancy"
% (running text), "plain" (front matter / chapter-openers / table of contents) and
% "empty" (title page + blank inter-chapter pages) -- so the footer also appears on
% the title page and the front matter.
%
% NOTE: macro names are deliberately @-free so they do not require \makeatletter
% (pandoc's include-in-header runs with a plain catcode for '@'). To change the
% footer wording, edit \LicenseFooter below and rebuild.

\usepackage{fancyhdr}

% The single source of truth for the footer wording.
\newcommand{\LicenseFooter}{%
      \small CC BY\ 4.0\ \textbullet\ \copyright\ 2026\ \ \textbf{Robert\ Ioffe}}

% One centered footer: license flush-left, page number flush-right, no rules.
\newcommand{\LicenseFancyStyle}{%
      \fancyhf{}%
      \fancyfoot[C]{\LicenseFooter \hfill \thepage}%
      \renewcommand{\headrulewidth}{0pt}}

% Apply the footer to every page style a book document uses.
\pagestyle{fancy}
\fancypagestyle{fancy}{\LicenseFancyStyle}
\fancypagestyle{plain}{\LicenseFancyStyle}
\fancypagestyle{empty}{\LicenseFancyStyle}
